% !TEX root = ../../main.tex

\section{Interpretation, Limits, and Next Checks}

The method has a clear inferential structure: build a hierarchy, annotate it
with child-parent and sibling tests, control for multiple testing, and convert
the annotated tree into a final partition. What remains essential is empirical
characterization of how this method behaves in practice.

Several questions matter for the paper. Under the null, the edge and
sibling tests should control false discoveries not only at the test level but
also at the level of the final number of clusters. Under planted structure,
the method should recover true splits often enough to be useful, especially as
sample size, signal strength, sparsity, and tree imbalance vary. On real and
synthetic benchmarks, the method should also show where it tends to over-split,
under-split, or lose power.

These empirical checks are particularly important here because the method is
not a single test applied in isolation. Tree construction, local testing,
multiple-testing correction, and top-down traversal interact. A convincing
evaluation therefore needs to describe the operating behavior of the full
method rather than only the calibration of one component.

The most informative empirical additions are null calibration, power under
known structure, robustness analyses under sparse and imbalanced regimes, and
benchmark comparisons tied directly to the exact implementation summarized in
Table~\ref{tab:defaults}. Together, these studies would show where the method
is well calibrated, where it remains powerful, and where its practical limits
begin.
